\subsection{Justificación del Problema}

% En este apartado debe explicar el por qué es necesario desarrollar este Trabajo de
% Titulación y/o la obtención del producto planteado, cuáles son los beneficios directos,
% la importancia teórica del proyecto, y otras consideraciones que crea necesario incluir.
% Es preciso enunciar los beneficios o impactos positivos de la investigación en el
% bienestar empresarial, regional, personal, conviene también añadir cifras, estadísticas
% y datos concretos que den validez a la justificación. Se explica además en que consiste
% el producto que se está planteando generar y los beneficios que este producto aportaría,
% pero, evitando entrar en detalles metodológicos u operacionales. En una justificación
% se anotará qué nuevos conocimientos aporta la investigación a la solución del problema

% planteado (importancia teórico-práctica), el significado que tiene la investigación para
% los beneficiarios, y quienes son estos, el por qué es útil la investigación (utilidad
% práctica), la factibilidad de desarrollar la investigación, es decir, aclarar si hay
% disponibilidad de recursos humanos, económicos, tecnológicos y bibliográficos
% actualizados. Se sugiere para la redacción basarse en las siguientes guías y redactar un
% párrafo de cada una de ellas:
% −Interés por investigar
% −Importancia teórico-práctica
% −Novedad en algún aspecto
% −Utilidad
% −Impacto
% - Factibilidad
% - Beneficiarios

% 1. Relevancia: por qué el problema merece estudiarse

% 2. Magnitud: a quién y cuánto afecta


% 3. Utilidad práctica: qué mejora concreta se espera


% 4. Aporte: qué conocimiento o método agrega


% 5. Factibilidad: por qué sí puede ejecutarse

%La digitalización de los procesos clínicos es hoy una prioridad reconocida en todo el mundo. Distintos organismos globales de salud destacan que tener una estrategia firme de salud digital mejora muchísimo los resultados médicos \cite{WHO2021DigitalHealth}. En el área de la odontología, la implementación de sistemas electrónicos para llevar los registros ha demostrado que revoluciona la gestión de los datos y la forma de atender, ya que facilita tener la documentación completa, aumenta la eficiencia del trabajo clínico y mejora la seguridad de las personas \cite{Sravya2024EHRDentistry}. Estudios recientes calculan que automatizar la historia clínica puede reducir entre un 35\% y un 40\% el tiempo que se gasta en tareas de oficina por cada paciente \cite{Sravya2024EHRDentistry}. Estos datos confirman que desarrollar un sistema informático integrado para Bellidental responde a una tendencia mundial y promete mejoras operativas muy importantes.

%El modelo actual de gestión manual en Bellidental causa problemas y trabas en el día a día que justifican por completo este estudio. La literatura médica muestra que usar historiales electrónicos en lugar de carpetas de papel incrementa la seguridad de los pacientes y la eficiencia del personal del consultorio. Por ejemplo, un estudio comparativo descubrió que las personas que usaron un sistema digital identificaron problemas de salud críticos mucho más rápido y sin olvidarse de nada, mientras que con los registros en papel cometieron un promedio de 9.3 omisiones en la información importante \cite{pmc2024digital}. Por el contrario, el flujo de trabajo manual donde el personal debe llevar los datos de un lado a otro causa demoras, tareas duplicadas y el riesgo constante de cometer errores. Además, mantener procesos separados y en papel va en contra de las recomendaciones internacionales sobre la protección de datos y las prácticas seguras de archivo. Todo esto demuestra que el problema actual de la clínica no es un detalle simple, sino algo que afecta la calidad del servicio y que merece una solución bien estructurada.

%El impacto de este problema afecta tanto al personal de la clínica como a las personas que van a atenderse. A nivel de los usuarios, los estudios de mercado revelan que hasta un 30\% de los pacientes de odontología pensaría en cambiarse de consultorio si se topa con retrasos frecuentes \cite{Minh2026WaitTimeDental}. De igual manera, asegurar que las citas empiecen a tiempo reduce en un 40\% las malas reseñas en internet que se quejan por las salas de espera llenas \cite{Minh2026WaitTimeDental}. En Bellidental, como todo depende de registros que están divididos en cuadernos, Excel y papel, es muy probable que existan demoras parecidas. Además, en los entornos de salud de países en desarrollo como el nuestro, el uso de sistemas integrados todavía es muy bajo, y se calcula que apenas el 40\% de los programas informáticos dentales pueden conectarse de verdad con otras historias clínicas \cite{Sravya2024EHRDentistry}. Esto significa que, si no se hace nada, los pacientes de Bellidental seguirán perdiendo tiempo en trámites y el consultorio tendrá problemas para hacerles un seguimiento correcto. Por lo tanto, este problema frena la satisfacción de los clientes, satura de trabajo al personal y limita el crecimiento de la clínica.

%Instalar un sistema de información clínico-administrativa centralizado traerá mejoras reales en el trabajo de todos los días. Si miramos casos reales en otros consultorios, con esta herramienta se espera bajar de forma notable el tiempo promedio que se pierde en el papeleo, logrando disminuciones de entre el 35\% y el 40\% en los trámites administrativos \cite{pmc2024digital}. Esto le dará al personal más libertad para concentrarse en la atención directa y el trato humano con los pacientes. Al mismo tiempo, ordenar los datos bajo un solo formato evita que se dupliquen archivos o se olviden datos médicos importantes. En las pruebas de los investigadores, el grupo que usó la historia clínica digital tuvo cero olvidos, frente a las 9.3 fallas de promedio del grupo que usó papel \cite{pmc2024digital}, lo que demuestra un avance claro en la seguridad médica. En resumen, el software integrado hará que la información se mueva de forma continua y transparente, eliminando los cuellos de botella para tener una atención rápida y registros mucho más confiables.

%Este proyecto aporta nuevos conocimientos y métodos que sirven tanto para la teoría como para la práctica. El hecho de aplicar un enfoque de gestión de procesos de negocio al caso de un consultorio dental pequeño es algo innovador en el contexto local. Los estudios previos demuestran que diseñar un sistema basado en procesos disminuye la necesidad de meter datos a mano, mejora la circulación de la información y optimiza todo el recorrido del paciente desde que llega hasta que se va \cite{makerere2024business}. Nuestro trabajo servirá para validar este enfoque directamente en la realidad de Bellidental. Además, la investigación ordenará las reglas de negocio y los requisitos de las leyes ecuatorianas, llenando un vacío importante, ya que se ha detectado que la literatura sobre software dental todavía está muy dispersa y le hace falta un análisis global \cite{Sravya2024EHRDentistry}. Por eso, la tesis creará un modelo de información que otros consultorios parecidos podrán copiar en el futuro.

%El proyecto es totalmente viable porque se cuenta con los recursos necesarios y existen ejemplos exitosos en la industria. Hoy en día hay herramientas tecnológicas muy maduras para armar estos sistemas, y se ha comprobado que usar plataformas en la nube reduce bastante los costos de computadoras y servidores para los negocios pequeños, lo que ayuda a la economía del proyecto \cite{Sravya2024EHRDentistry}. También se cuenta con bases académicas claras sobre cómo analizar los flujos de trabajo en salud que guiarán los primeros meses de la investigación \cite{makerere2024business}. Por la parte de las leyes, el Ministerio de Salud de Ecuador cuenta con un marco legal claro para el manejo de información confidencial, lo que asegura que el diseño del programa respetará la privacidad desde el primer día. Finalmente, el equipo de Bellidental está totalmente comprometido a colaborar con las entrevistas y las pruebas del software, lo que garantiza que el desarrollo de la herramienta es factible en el tiempo planificado.

%En conclusión, la razón de ser de este trabajo está en los beneficios reales que va a dejar de forma inmediata, mejorando la eficiencia interna y la calidad de la atención médica en el consultorio. La evidencia de las investigaciones respalda que digitalizar los procesos administrativos y clínicos disminuye el papeleo pesado y cuida mejor al paciente \cite{Sravya2024EHRDentistry}. Al diseñar este sistema integrado para Bellidental, se logrará un servicio más ágil, con menos esperas, libre de equivocaciones y alineado con las leyes de salud, permitiendo que el centro crezca sin las trabas que tiene actualmente.

La digitalización de los procesos clínicos es hoy una prioridad reconocida
en todo el mundo. Distintos organismos globales de salud destacan que
contar con una estrategia firme de salud digital mejora de manera
significativa los resultados médicos \cite{WHO2021DigitalHealth}. En el
área de la odontología, la implementación de sistemas electrónicos de
registro ha demostrado revolucionar la gestión de datos y la forma de
atender, facilitando la documentación completa, incrementando la eficiencia
del trabajo clínico y mejorando la seguridad del paciente
\cite{Sravya2024EHRDentistry}. La literatura especializada estima, como
referencia teórica, que automatizar la historia clínica tiene el potencial
de reducir entre un 35\,\% y un 40\,\% el tiempo destinado a tareas
administrativas por paciente \cite{Sravya2024EHRDentistry}; esta cifra
opera como horizonte de mejora que orienta el presente proyecto, cuya
magnitud real de reducción será determinada empíricamente mediante la
medición de los minutos por ciclo completo de atención antes y después
de la implementación del sistema.

El modelo actual de gestión manual en Bellidental genera ineficiencias
concretas que justifican esta investigación. La literatura médica demuestra
que el uso de registros electrónicos en lugar de carpetas de papel
incrementa la seguridad de los pacientes y la eficiencia del personal.
Un estudio comparativo mostró que los usuarios de sistemas digitales
identificaron condiciones de salud críticas con mayor rapidez y sin
omisiones, mientras que quienes trabajaron con registros en papel
registraron un promedio de 9,3 omisiones de información relevante por
paciente \cite{pmc2024digital}. Por el contrario, el flujo manual —donde
el personal traslada datos entre herramientas desconectadas— genera
demoras, tareas duplicadas y riesgo constante de error. Mantener procesos
dispersos en papel es, además, incompatible con las recomendaciones
internacionales de protección de datos y archivo seguro. Estas evidencias
confirman que el problema de Bellidental afecta directamente la calidad
del servicio y demanda una solución estructurada.

El impacto de esta fragmentación recae tanto en el personal como en los
pacientes. Estudios de mercado en odontología revelan que hasta un 30\,\%
de los pacientes consideraría cambiar de consultorio ante retrasos
frecuentes, y que garantizar la puntualidad en las citas reduce en un
40\,\% las reseñas negativas asociadas a tiempos de espera excesivos
\cite{Minh2026WaitTimeDental}. En Bellidental, donde el registro de un
paciente exige buscar manualmente entre miles de filas de Excel, coordinar
citas en Google Calendar y anotar pagos en papel, es previsible que este
fenómeno ya esté ocurriendo. A escala regional, se estima que apenas el
40\,\% de los sistemas informáticos dentales en entornos de desarrollo
logra conectarse con otros repositorios de historias clínicas
\cite{Sravya2024EHRDentistry}, lo que confirma que la fragmentación no
es un problema aislado, sino estructural. El indicador que sintetiza todas
estas consecuencias es el tiempo que el personal invierte por cada
paciente en el ciclo completo de atención: reducirlo es reducir
simultáneamente las esperas, los errores y la sobrecarga operativa.

La implementación de un sistema centralizado producirá mejoras medibles
en la operación diaria. Consolidar las historias clínicas, las citas, las
recetas y los cobros en una sola plataforma elimina la necesidad de
trasladar datos manualmente entre herramientas, lo que acorta directamente
el tiempo de cada ciclo de atención. En estudios controlados, el grupo
con historia clínica digital registró cero omisiones frente a las 9,3
del grupo con papel \cite{pmc2024digital}, evidenciando ganancias
simultáneas en velocidad y en calidad del registro. Al reducir los minutos
invertidos en papeleo por paciente, el personal recupera tiempo para
la atención directa, y la clínica gana capacidad para absorber mayor
demanda sin ampliar su planta.

Este proyecto aporta valor teórico y práctico en el contexto local.
Aplicar un enfoque de gestión de procesos de negocio a un consultorio
dental de pequeña escala es una contribución novedosa para la región,
dado que los estudios previos confirman que modelar el flujo real del
paciente antes de implementar tecnología reduce los desajustes entre el
sistema y la operación, y mejora la aceptación del personal
\cite{makerere2024business}. Adicionalmente, la investigación formalizará
las reglas de negocio y los requisitos normativos del contexto ecuatoriano,
llenando un vacío identificado en la literatura, que señala la dispersión
y falta de síntesis global en el campo de la informática odontológica
\cite{Sravya2024EHRDentistry}. El modelo de información resultante podrá
ser replicado por consultorios de características similares.

La viabilidad del proyecto está respaldada por factores técnicos,
normativos y organizacionales. Existen herramientas tecnológicas maduras
para el desarrollo de estos sistemas, y el uso de plataformas en la nube
permite reducir costos de infraestructura para negocios pequeños
\cite{Sravya2024EHRDentistry}. El marco metodológico para el análisis de
flujos de trabajo en salud está bien documentado en la literatura académica
\cite{makerere2024business}, y el Ministerio de Salud Pública del Ecuador
provee un marco normativo claro —el Acuerdo Ministerial 5216-A y el
Formulario 033— que orienta el diseño desde la primera fase. El equipo
de Bellidental ha manifestado disponibilidad total para participar en las
entrevistas, observaciones y pruebas del sistema, lo que hace factible
cumplir el cronograma planificado.

En conclusión, este trabajo se justifica por los beneficios directos y
cuantificables que generará en la operación del consultorio. La evidencia
disponible respalda que centralizar los procesos clínico-administrativos
en un sistema integrado reduce la carga documental y mejora la seguridad
del paciente \cite{Sravya2024EHRDentistry}. El efecto de esa centralización
será medido a través de la variable que concentra todas las consecuencias
del problema actual: los minutos por ciclo completo de atención por
paciente. Reducir esa cifra es el objetivo concreto y verificable que
este proyecto se propone demostrar.